\documentclass{beamer}

\usepackage[english]{babel}
\usepackage[T1]{fontenc}
\usepackage{graphicx}
\usepackage{ifthen}
\usepackage{listings}
\usepackage{verbatim}
\usepackage{fancyhdr}
\usepackage{tikz}
\usepackage{dsfont}
\usetikzlibrary{spy,matrix}
\usepackage[style=mla, backend=bibtex]{biblatex}
\bibliography{storm}

\usepackage{xunicode}
\usepackage{xltxtra}
\usepackage{setspace}
\defaultfontfeatures{Mapping=tex-text}
\setmonofont[Mapping={}, Scale=MatchLowercase]{Source Code Pro}
\setsansfont[Scale=MatchLowercase]{Source Sans Pro}
\setmainfont[]{Source Serif Pro}

\newbox\mytempbox
\newdimen\mytempdimen
\newcommand\includegraphicscopyright[3][]{%
  \leavevmode\vbox{\vskip3pt\raggedright\setbox\mytempbox=\hbox{%
  \includegraphics[#1]{#2}}%
    \mytempdimen=\wd\mytempbox\box\mytempbox\par\vskip1pt%
    \fontsize{3}{3.5}\selectfont{\color{black!25}{\vbox{\hsize=\mytempdimen#3}}}\vskip3pt%
}}

\newcommand\prelim[1]{\small }

\newcommand\strColor[1]{\color{beamer@blendedblue}{#1}}

\newcommand\sect[1]{\begin{center}\huge\strColor{#1}\end{center}}

\setbeamertemplate{navigation symbols}{}

\pagestyle{fancy}
\rhead{\vspace{.3em}\includegraphics[width=10em]{assets/FULogo_RGB.eps}}
\cfoot{}
\renewcommand{\headrulewidth}{0pt}

\lstset{language=C}
\lstset{basicstyle=\footnotesize\ttfamily}
\lstset{breaklines=true}
\lstset{keywordstyle=\color{purple}}

\usepackage{geometry}
\usepackage{tikz}
\usetikzlibrary{calc,arrows}

\usepackage{pifont}

\newdimen\XCoord
\newdimen\YCoord

\newcommand*{\ExtractCoordinate}[1]{\path (#1); \pgfgetlastxy{\XCoord}{\YCoord};}%

\title{POSIX compliant shell}
\author{Robert Fels, Martine Lenders, Jakob Pfender}
\institute{FU Berlin\\Institut für Informatik\\Softwareprojekt
Telematik}
\date{\today}

\begin{document}

\frame{\titlepage}

\frame{
\frametitle{Project topic}
\begin{itemize}
  \item \textbf{Goal:} Port a POSIX compliant shell (dash) to RIOT
  \begin{itemize}
    \item \textbf{Motivation:} provide an execution environment for shell scripts
    \item \textbf{Motivation:} Improve RIOT's POSIX compliance
  \end{itemize}
\end{itemize}
}

\frame{
    \frametitle{Work schedule -- Three phases}
    \begin{itemize}
      \item \textbf{Phase 1:} Compilation
        \begin{itemize}
          \item \textbf{Subphase 1:} Identify missing headers $\checkmark$
          \item \textbf{Subphase 2:} Implement missing headers ($\checkmark$) (prototypes)
        \end{itemize}
      \item \textbf{Phase 2:} Functionality
        \begin{itemize}
          \item \textbf{Subphase 3:} Identify missing functionality
            $\checkmark$
          \item \textbf{Subphase 4:} Implement missing functionality ($\checkmark$)
            \begin{itemize}
              \item Process handling and signaling $\checkmark$
                \begin{itemize}
                  \item \texttt{fork()} $\checkmark$
                  \item \texttt{kill()} $\checkmark$
                \end{itemize}
              \item File system ($\checkmark$)
              \item Way to execute RIOT's predefined shell handlers
                (\texttt{shell\_commands}) ($\checkmark$)
              \item Wrappers for POSIX functions $\times$
            \end{itemize}
          \item \textbf{Subphase 5:} Test $\times$
        \end{itemize}
    \end{itemize}
}

\frame{
  \frametitle{Issues}
  \begin{itemize}
    \item Close to none of the POSIX prerequisites needed to implement
      a POSIX compliant shell were there to begin with
    \item RIOT process handling misses a lot of the functionality
      expected by POSIX
      \begin{itemize}
        \item e.g. no signaling, no forking
      \end{itemize}
    \item No file system (UNIX philosophy: everything is a file)
  \end{itemize}
}

\frame{
  \frametitle{Achievements (1)}
  \begin{itemize}
    \item Defined POSIX prototypes
    \item Implemented capabilities to kill processes
        \begin{itemize}
            \item Pull request submitted
            \item stalled, since IPC can't handle threads getting suddenly killed
            \item still waiting for PR \#1270
        \end{itemize}
    \item Forking
      \begin{itemize}
        \item Pull request submitted
      \end{itemize}
  \end{itemize}
}

\frame{
  \frametitle{Achievements (2)}
  \begin{itemize}
    \item Compilation of *.c and *.h files
    \item Designed file system interface (loosely based on FUSE)
    \item informed about possible file systems to implement to present shell
          scripts to RIOT
          \begin{itemize}
              \item FAT (problematic because of licensing issues)
              \item UFFS, nuttX (not used so often, quite new, missing thorough testing)
	      \item Coffee file system (CFS) from contiki seems to a good choice
	      \begin{itemize}
		\item minimalistic, yet fully functional file system
		\item POSIX complient
		\item primarily aimed at sensor devices that equiped with Flash Memories or EEPROM
		\item there are no Contiki-specific abstractions used in the implementation 			      of CFS
		\item implementation not stated yet
	      \end{itemize}
          \end{itemize}
    \item SCFS (shell\_commands file system)
      \begin{itemize}
        \item Pull request submitted
      \end{itemize}
  \end{itemize}
}

\frame{
  \frametitle{What was planned v. what was implemented}
  \begin{itemize}
    \item No wrappers for POSIX functionality
    \item No functioning file system
  \end{itemize}
}

\frame{
  \frametitle{Perspectives}
  \begin{itemize}
    \item Keep working on currently open pull requests
      \begin{itemize}
        \item Refactoring
        \item Documentation
        \item Tests
      \end{itemize}
    \item Implement file system and missing POSIX functionalities
  \end{itemize}
}

\frame{
  \frametitle{Pull Requests}
  \begin{itemize}
    \item \textbf{(Include missing functions in header files):} \#1166
    \item \textbf{Return-based error handling for shell handlers:} \#1243
    \item \textbf{Killing processes:} \#1245
    \item \textbf{File system:} \#1265
    \item \textbf{Fork:} \#1384
    \item \textbf{SCFS:} \#1418 (depends on \#1265)
  \end{itemize}
}

\frame{
  \huge{Questions?}
}

\end{document}
